\documentclass[11pt]{article}

\usepackage{amsmath, amssymb, amsthm}
\usepackage{mathtools}
\usepackage{geometry}
\usepackage{enumitem}

\geometry{margin=1in}

\newtheorem{theorem}{Theorem}[section]
\newtheorem{proposition}[theorem]{Proposition}
\newtheorem{lemma}[theorem]{Lemma}
\newtheorem{corollary}[theorem]{Corollary}
\theoremstyle{definition}
\newtheorem{definition}[theorem]{Definition}
\newtheorem{example}[theorem]{Example}
\newtheorem{remark}[theorem]{Remark}

\DeclareMathOperator{\Tr}{Tr}
\DeclareMathOperator{\Range}{Range}

\newcommand{\R}{\mathbb{R}}
\newcommand{\N}{\mathbb{N}}
\newcommand{\E}{\mathbb{E}}
\newcommand{\Splus}{\mathcal{S}_+^d}
\newcommand{\Splusplus}{\mathcal{S}_{++}^d}
\newcommand{\HK}{\mathcal{H}_K}
\newcommand{\X}{\mathcal{X}}
\newcommand{\B}{\mathcal{B}}
\newcommand{\M}{\mathcal{M}}
\newcommand{\iid}{\stackrel{\mathrm{iid}}{\sim}}
\newcommand{\inprod}[2]{\langle #1, #2 \rangle}
\newcommand{\norm}[1]{\| #1 \|}

\title{Integral Anisotropic Gaussian Kernels:\\Formal Definitions}
\author{}
\date{}

\begin{document}

\maketitle

%%%%%%%%%%%%%%%%%%%%%%%%%%%%%%%%%%%%%%%%%%%%%%%%%%%%%%%%%%%%%%%%%%%%%%%%%%%%%%%
\section{Preliminaries}
%%%%%%%%%%%%%%%%%%%%%%%%%%%%%%%%%%%%%%%%%%%%%%%%%%%%%%%%%%%%%%%%%%%%%%%%%%%%%%%

\begin{definition}[Cone of Positive Semi-Definite Matrices]
Let $d \geq 1$. The \emph{cone of symmetric positive semi-definite matrices} is
\begin{equation}
    \Splus = \left\{ A \in \R^{d \times d} : A = A^\top \text{ and } x^\top A x \geq 0 \;\; \forall x \in \R^d \right\}.
\end{equation}
The interior $\Splusplus$ denotes the \emph{positive definite} matrices (strict inequality for $x \neq 0$).
\end{definition}

\begin{definition}[Anisotropic Gaussian Kernel]
For $A \in \Splus$, the \emph{anisotropic Gaussian kernel} with precision matrix $A$ is
\begin{equation}
    K_A: \R^d \times \R^d \to \R, \quad K_A(x, y) = \exp\left( -\frac{1}{2}(x-y)^\top A (x-y) \right).
\end{equation}
\end{definition}

\begin{definition}[Borel Measures on $\Splus$]
Equip $\Splus$ with the subspace topology inherited from $\R^{d \times d}$ and let $\B(\Splus)$ denote the Borel $\sigma$-algebra. Let $\M_+(\Splus)$ denote the space of finite non-negative Borel measures on $\Splus$, and let $\mathcal{P}(\Splus) \subset \M_+(\Splus)$ denote the probability measures.
\end{definition}

%%%%%%%%%%%%%%%%%%%%%%%%%%%%%%%%%%%%%%%%%%%%%%%%%%%%%%%%%%%%%%%%%%%%%%%%%%%%%%%
\section{The Integral Anisotropic Gaussian Kernel}
%%%%%%%%%%%%%%%%%%%%%%%%%%%%%%%%%%%%%%%%%%%%%%%%%%%%%%%%%%%%%%%%%%%%%%%%%%%%%%%

\begin{definition}[Integral Anisotropic Gaussian Kernel]
\label{def:integral-kernel}
Let $\theta \in \M_+(\Splus)$ be a finite non-negative Borel measure on the cone of positive semi-definite matrices. The \emph{integral anisotropic Gaussian kernel} induced by $\theta$ is the function $K_\theta: \R^d \times \R^d \to \R$ defined by
\begin{equation}
    \boxed{
    K_\theta(x, y) = \int_{\Splus} K_A(x, y) \, d\theta(A) = \int_{\Splus} \exp\left( -\frac{1}{2}(x-y)^\top A (x-y) \right) d\theta(A).
    }
\end{equation}
We call $\theta$ the \emph{mixing measure} or \emph{structure measure} of the kernel.
\end{definition}

\begin{remark}
Special cases:
\begin{enumerate}[label=(\roman*)]
    \item If $\theta = \sum_{j=1}^J a_j \delta_{A_j}$ is a finite weighted sum of Dirac masses, then
    \begin{equation}
        K_\theta(x,y) = \sum_{j=1}^J a_j \, K_{A_j}(x,y),
    \end{equation}
    recovering the finite mixture of anisotropic Gaussian kernels.
    \item If $\theta \in \mathcal{P}(\Splus)$ is a probability measure, $K_\theta$ is a \emph{scale-mixture} of anisotropic Gaussians.
\end{enumerate}
\end{remark}

%%%%%%%%%%%%%%%%%%%%%%%%%%%%%%%%%%%%%%%%%%%%%%%%%%%%%%%%%%%%%%%%%%%%%%%%%%%%%%%
\section{Well-Definedness and Basic Properties}
%%%%%%%%%%%%%%%%%%%%%%%%%%%%%%%%%%%%%%%%%%%%%%%%%%%%%%%%%%%%%%%%%%%%%%%%%%%%%%%

\begin{proposition}[Well-Definedness]
\label{prop:well-defined}
For any $\theta \in \M_+(\Splus)$, the integral in Definition~\ref{def:integral-kernel} is well-defined and finite for all $x, y \in \R^d$.
\end{proposition}

\begin{proof}
For all $A \in \Splus$ and $x, y \in \R^d$, we have $(x-y)^\top A (x-y) \geq 0$, hence
\begin{equation}
    0 < K_A(x,y) = \exp\left(-\frac{1}{2}(x-y)^\top A (x-y)\right) \leq 1.
\end{equation}
Therefore,
\begin{equation}
    0 < K_\theta(x,y) \leq \int_{\Splus} 1 \, d\theta(A) = \theta(\Splus) < \infty. \qedhere
\end{equation}
\end{proof}

\begin{proposition}[Positive Semi-Definiteness]
\label{prop:psd}
For any $\theta \in \M_+(\Splus)$, the kernel $K_\theta$ is positive semi-definite.
\end{proposition}

\begin{proof}
Each $K_A$ is positive semi-definite. For any $n \in \N$, points $\{x_i\}_{i=1}^n \subset \R^d$, and coefficients $\{c_i\}_{i=1}^n \subset \R$:
\begin{equation}
    \sum_{i,j=1}^n c_i c_j K_\theta(x_i, x_j) 
    = \int_{\Splus} \underbrace{\sum_{i,j=1}^n c_i c_j K_A(x_i, x_j)}_{\geq 0 \text{ since } K_A \text{ is PSD}} d\theta(A) \geq 0. \qedhere
\end{equation}
\end{proof}

\begin{proposition}[Stationarity]
\label{prop:stationary}
The kernel $K_\theta$ is stationary: $K_\theta(x, y) = \phi_\theta(x - y)$ where
\begin{equation}
    \phi_\theta(z) = \int_{\Splus} \exp\left(-\frac{1}{2} z^\top A z\right) d\theta(A).
\end{equation}
\end{proposition}

\begin{proposition}[Normalization]
$K_\theta(x, x) = \theta(\Splus)$ for all $x \in \R^d$.
\end{proposition}

%%%%%%%%%%%%%%%%%%%%%%%%%%%%%%%%%%%%%%%%%%%%%%%%%%%%%%%%%%%%%%%%%%%%%%%%%%%%%%%
\section{Fourier Representation}
%%%%%%%%%%%%%%%%%%%%%%%%%%%%%%%%%%%%%%%%%%%%%%%%%%%%%%%%%%%%%%%%%%%%%%%%%%%%%%%

\begin{proposition}[Spectral Representation]
\label{prop:fourier}
If $\theta$ is supported on $\Splusplus$ and satisfies $\int_{\Splusplus} |A|^{-1/2} d\theta(A) < \infty$, then
\begin{equation}
    K_\theta(x, y) = \int_{\R^d} e^{i\omega^\top(x-y)} \, p_\theta(\omega) \, d\omega,
\end{equation}
where the \emph{spectral density} is the mixture of Gaussians:
\begin{equation}
    p_\theta(\omega) = \int_{\Splusplus} \frac{|A|^{1/2}}{(2\pi)^{d/2}} \exp\left(-\frac{1}{2}\omega^\top A^{-1} \omega\right) d\theta(A).
\end{equation}
\end{proposition}

%%%%%%%%%%%%%%%%%%%%%%%%%%%%%%%%%%%%%%%%%%%%%%%%%%%%%%%%%%%%%%%%%%%%%%%%%%%%%%%
\section{The RKHS of $K_\theta$}
%%%%%%%%%%%%%%%%%%%%%%%%%%%%%%%%%%%%%%%%%%%%%%%%%%%%%%%%%%%%%%%%%%%%%%%%%%%%%%%

\begin{theorem}[RKHS Structure]
\label{thm:rkhs}
The reproducing kernel Hilbert space $\mathcal{H}_{K_\theta}$ of $K_\theta$ admits the characterization:
\begin{equation}
    \mathcal{H}_{K_\theta} = \left\{ f = \int_{\Splus} f_A \, d\theta(A) : f_A \in \mathcal{H}_{K_A} \text{ for } \theta\text{-a.e. } A, \;\; \int_{\Splus} \norm{f_A}_{\mathcal{H}_{K_A}}^2 d\theta(A) < \infty \right\}
\end{equation}
with norm given by the \emph{continuous infimal convolution}:
\begin{equation}
    \boxed{
    \norm{f}_{\mathcal{H}_{K_\theta}}^2 = \inf \left\{ \int_{\Splus} \norm{f_A}_{\mathcal{H}_{K_A}}^2 \, d\theta(A) : f = \int_{\Splus} f_A \, d\theta(A), \;\; f_A \in \mathcal{H}_{K_A} \right\}.
    }
\end{equation}
\end{theorem}

%%%%%%%%%%%%%%%%%%%%%%%%%%%%%%%%%%%%%%%%%%%%%%%%%%%%%%%%%%%%%%%%%%%%%%%%%%%%%%%
\section{The Associated Gaussian Measure}
%%%%%%%%%%%%%%%%%%%%%%%%%%%%%%%%%%%%%%%%%%%%%%%%%%%%%%%%%%%%%%%%%%%%%%%%%%%%%%%

\begin{definition}[Gaussian Measure Induced by $K_\theta$]
\label{def:gp}
The \emph{Gaussian measure} $\mu_{K_\theta}$ associated with $K_\theta$ is the unique Borel probability measure on $L^2(P_X)$ with characteristic functional
\begin{equation}
    \hat{\mu}_{K_\theta}(g) = \int_{L^2} e^{i\inprod{g}{f}_{L^2}} d\mu_{K_\theta}(f) = \exp\left(-\frac{1}{2}\inprod{g}{T_{K_\theta} g}_{L^2}\right),
\end{equation}
where $T_{K_\theta}$ is the integral operator with kernel $K_\theta$.
\end{definition}

\begin{theorem}[Decomposition of the Gaussian Process]
\label{thm:gp-decomposition}
If $\theta \in \mathcal{P}(\Splus)$ is a probability measure, then
\begin{equation}
    f \sim \mu_{K_\theta} \iff f = \int_{\Splus} f_A \, d\theta(A), \quad \text{where } \{f_A\}_{A \in \Splus} \text{ are independent with } f_A \sim \mu_{K_A}.
\end{equation}
More precisely, there exists a probability space supporting a random field $(f_A)_{A \in \Splus}$ such that:
\begin{enumerate}[label=(\roman*)]
    \item For $\theta$-almost every $A$, $f_A \sim \mathcal{GP}(0, K_A)$.
    \item The marginals $(f_{A_1}, \ldots, f_{A_n})$ are independent for $\theta$-a.e.\ distinct $A_1, \ldots, A_n$.
    \item $f = \int_{\Splus} f_A \, d\theta(A)$ satisfies $f \sim \mathcal{GP}(0, K_\theta)$.
\end{enumerate}
\end{theorem}

%%%%%%%%%%%%%%%%%%%%%%%%%%%%%%%%%%%%%%%%%%%%%%%%%%%%%%%%%%%%%%%%%%%%%%%%%%%%%%%
\section{Common Choices of Mixing Measure $\theta$}
%%%%%%%%%%%%%%%%%%%%%%%%%%%%%%%%%%%%%%%%%%%%%%%%%%%%%%%%%%%%%%%%%%%%%%%%%%%%%%%

\begin{example}[Finite Mixture]
\begin{equation}
    \theta = \sum_{j=1}^J a_j \delta_{A_j}, \quad a_j > 0, \quad A_j \in \Splusplus.
\end{equation}
\end{example}

\begin{example}[Isotropic Scale Mixture]
Let $\nu$ be a measure on $(0, \infty)$ and set $\theta = \nu \circ \psi^{-1}$ where $\psi(s) = s \cdot I_d$. Then:
\begin{equation}
    K_\theta(x,y) = \int_0^\infty \exp\left(-\frac{s}{2}\norm{x-y}^2\right) d\nu(s).
\end{equation}
This recovers Schoenberg's representation for radial positive definite functions.
\end{example}

\begin{example}[Wishart Mixing Measure]
Let $\theta = \text{Wishart}_d(\nu, \Psi)$ with density
\begin{equation}
    p(A) = \frac{|A|^{(\nu-d-1)/2}}{2^{\nu d/2} |\Psi|^{\nu/2} \Gamma_d(\nu/2)} \exp\left(-\frac{1}{2}\Tr(\Psi^{-1} A)\right), \quad A \in \Splusplus,
\end{equation}
where $\nu > d - 1$ and $\Gamma_d$ is the multivariate gamma function.
\end{example}

\begin{example}[Rank-One (Ridge Function) Measure]
Let $\lambda > 0$ and let $\rho$ be a measure on the sphere $\mathbb{S}^{d-1}$. Define $\theta$ as the pushforward of $\rho$ under $w \mapsto \lambda \cdot ww^\top$:
\begin{equation}
    K_\theta(x,y) = \int_{\mathbb{S}^{d-1}} \exp\left(-\frac{\lambda}{2}(w^\top(x-y))^2\right) d\rho(w).
\end{equation}
This connects to neural networks and ridge function approximation (Bach, 2017).
\end{example}

\begin{example}[Spectral Decomposition Measure]
Write $A = U \Lambda U^\top$ with $U \in O(d)$ orthogonal and $\Lambda = \text{diag}(\lambda_1, \ldots, \lambda_d)$. Let
\begin{equation}
    d\theta(A) = d\mu(\lambda_1, \ldots, \lambda_d) \, d\nu(U),
\end{equation}
where $\mu$ is a measure on $\R_+^d$ and $\nu$ is a measure on the orthogonal group $O(d)$.

If $\nu$ is the Haar measure, then $K_\theta$ is \emph{rotation-invariant}.
\end{example}

%%%%%%%%%%%%%%%%%%%%%%%%%%%%%%%%%%%%%%%%%%%%%%%%%%%%%%%%%%%%%%%%%%%%%%%%%%%%%%%
\section{Characterization Theorem}
%%%%%%%%%%%%%%%%%%%%%%%%%%%%%%%%%%%%%%%%%%%%%%%%%%%%%%%%%%%%%%%%%%%%%%%%%%%%%%%

\begin{theorem}[Representation Theorem]
\label{thm:representation}
A continuous stationary kernel $K(x,y) = \phi(x-y)$ on $\R^d$ can be represented as
\begin{equation}
    \phi(z) = \int_{\Splus} \exp\left(-\frac{1}{2}z^\top A z\right) d\theta(A)
\end{equation}
for some $\theta \in \M_+(\Splus)$ if and only if its spectral measure (via Bochner's theorem) is a mixture of Gaussian measures on $\R^d$.

Equivalently, $\phi$ is representable if and only if its Fourier transform satisfies
\begin{equation}
    \hat{\phi}(\omega) = \int_{\Splusplus} \exp\left(-\frac{1}{2}\omega^\top \Sigma \omega\right) d\rho(\Sigma)
\end{equation}
for some finite measure $\rho$ on $\Splusplus$ (where $\Sigma = A^{-1}$).
\end{theorem}

\begin{corollary}[Radial Case: Schoenberg-Bernstein]
If $\phi(z) = \psi(\norm{z}^2)$ is radial, then $\phi$ is representable as a mixture of \emph{isotropic} Gaussians if and only if $\psi: [0,\infty) \to \R$ is completely monotone, i.e.,
\begin{equation}
    (-1)^n \psi^{(n)}(t) \geq 0 \quad \forall t > 0, \; \forall n \geq 0.
\end{equation}
\end{corollary}

\begin{theorem}[Density]
\label{thm:density}
The family $\{K_\theta : \theta \in \M_+(\Splus)\}$ is dense in the space of continuous stationary positive semi-definite kernels on $\R^d$, under the topology of uniform convergence on compact sets.
\end{theorem}

%%%%%%%%%%%%%%%%%%%%%%%%%%%%%%%%%%%%%%%%%%%%%%%%%%%%%%%%%%%%%%%%%%%%%%%%%%%%%%%
\section{Summary}
%%%%%%%%%%%%%%%%%%%%%%%%%%%%%%%%%%%%%%%%%%%%%%%%%%%%%%%%%%%%%%%%%%%%%%%%%%%%%%%

\begin{center}
\fbox{\parbox{0.92\textwidth}{
\textbf{Definition.} For $\theta \in \M_+(\Splus)$, the integral anisotropic Gaussian kernel is:
\begin{equation*}
    K_\theta(x, y) = \int_{\Splus} \exp\left(-\frac{1}{2}(x-y)^\top A (x-y)\right) d\theta(A).
\end{equation*}

\vspace{0.3em}
\textbf{RKHS norm} (continuous infimal convolution):
\begin{equation*}
    \norm{f}_{\mathcal{H}_{K_\theta}}^2 = \inf_{f = \int f_A \, d\theta(A)} \int_{\Splus} \norm{f_A}_{\mathcal{H}_{K_A}}^2 \, d\theta(A).
\end{equation*}

\vspace{0.3em}
\textbf{GP decomposition}: $f \sim \mathcal{GP}(0, K_\theta) \iff f = \int f_A \, d\theta(A)$ with $f_A \stackrel{\text{ind}}{\sim} \mathcal{GP}(0, K_A)$.

\vspace{0.3em}
\textbf{Spectral density}: $p_\theta(\omega) = \int_{\Splusplus} \frac{|A|^{1/2}}{(2\pi)^{d/2}} e^{-\frac{1}{2}\omega^\top A^{-1}\omega} d\theta(A)$.
}}
\end{center}

\end{document}
